\documentclass[11pt,a4paper]{article}
\usepackage[utf8]{inputenc}
\usepackage[T1]{fontenc}
\usepackage[margin=2.4cm]{geometry}
\usepackage{amsmath,amssymb,amsthm,mathtools,booktabs,array,xcolor}
\emergencystretch=2em
\usepackage[colorlinks=true,linkcolor=blue!60!black,citecolor=green!40!black,urlcolor=blue!60!black]{hyperref}

\newcommand{\La}{\textsf{\textbf{(a)}}}
\newcommand{\Lb}{\textsf{\textbf{(b)}}}
\newcommand{\Lc}{\textsf{\textbf{(c)}}}
\newcommand{\Ld}{\textsf{\textbf{(d)}}}
\newcommand{\dd}{\mathrm{d}}
\newcommand{\ii}{\mathrm{i}}
\newcommand{\ee}{\mathrm{e}}
\newcommand{\hb}{\hbar}
\newcommand{\sgn}{\operatorname{sgn}}
\newcommand{\Rea}{\operatorname{Re}}
\newcommand{\Ima}{\operatorname{Im}}
\newcommand{\script}[1]{\texttt{phase2/\allowbreak scripts/\allowbreak #1}}
\newcommand{\sclass}{\texttt{phase2/\allowbreak scripts/\allowbreak defq\_\allowbreak classical\_\allowbreak charges.py}}
\newcommand{\sren}{\texttt{phase2/\allowbreak scripts/\allowbreak defq\_\allowbreak oneloop\_\allowbreak renorm.py}}
\newcommand{\sfus}{\texttt{phase2/\allowbreak scripts/\allowbreak defq\_\allowbreak absorption\_\allowbreak fusion.py}}

\newtheorem{proposition}{Proposition}[section]
\theoremstyle{remark}
\newtheorem{remark}[proposition]{Remark}

\title{\textbf{Boundary Charges, One-Loop Renormalisation,\\
and Obstructions to Quantum Defect Fusion\\
for the Corrigan--Zambon Boussinesq Defect}}
\author{Project \emph{quantum\_boussinesq\_solitons} (Route E, next step \#3) --- PI: Gustav Delius}
\date{4 November 2026}

\begin{document}
\maketitle

\begin{abstract}
We carry out ranked next step~\#3 of the revised Phase~2 master report
(\texttt{phase2/\allowbreak report.pdf}): deriving the localised classical and one-loop
boundary charges and renormalisation structure of the Corrigan--Zambon (CZ)
type-I integrable defect \cite{CZ} for the good Boussinesq equation before
making any physical claim of quantum defect fusion. Every nontrivial claim is
labelled \La\ (literature), \Lb\ (literature heuristic), \Lc\ (derived and
verified in this project), or \Ld\ (guess/open physical identification).

Our main results \Lc\ are:
\begin{enumerate}
\item \textbf{Classical localised defect charges and exact conservation laws (\sclass).}
Starting from the total action in fixed project conventions ($u_{\mathrm{CZ}}=-3\varphi$,
$\mathcal{L}_{\mathrm{CZ}}=9\mathcal{L}$), we derive the four nonlinear sewing
conditions and the exact localised defect energy $E_D$, modified momentum $P_D$,
$\varphi$-shift charge $\Pi_D$, and topological mass jump $M_D = \varphi_R(0^+)-\varphi_L(0^-)$.
We prove symbolically that $\frac{\dd}{\dd t}(E_{\mathrm{bulk}}+E_D)=0$,
$\frac{\dd}{\dd t}(P_{\mathrm{bulk}}+P_D)=0$,
$\frac{\dd}{\dd t}(\Pi_{\mathrm{bulk}}+\Pi_D)=0$, and
$\frac{\dd}{\dd t}(M_{\mathrm{bulk}}+M_D)=0$ hold off-shell on the sewing
manifold, that the vacuum jumps across any KP triple edge $S_i=\{q_i,q_{i+1}\}$
factorise modulo $q_i^2+q_i q_{i+1}+q_{i+1}^2=\frac{1}{4}$ into
$\Delta(E_D,P_D,\Pi_D,M_D)_{q_i\to q_{i+1}} = \rho_i(E,P,\Pi,M)_{S_i}$ with
$\rho_i = \eta\sgn(q_i-q_{i+1})$, and we verify time-independent total charges
to $10^{-50}$ on exact transmitting and absorbing soliton--defect solutions.
\item \textbf{Corrected sewing determinant, admissible growing-mode branch switch, and exact spectral trace of $P_{D,2}$ (\sren).}
Incorporating the independent audit (\texttt{audit\_frames\_defects.pdf}), we
verify symbolically that the reduced $4\times 4$ linearised sewing determinant
carries the prefactor $-\frac{2}{3}$ (sixteen times the Phase~2g prefactor
$-\frac{1}{24}$) and that its continuum-threshold values factorise as
$P_-(1;s)=\frac{1}{2}(s+1)(s-\sqrt{3})$ and
$P_+(1;s)=\frac{1}{2}(s-1)(s-\sqrt{3})$ ($s=\eta\sigma_{\mathrm{vac}}$). This
proves an exact admissible-branch switch ($S=\kappa_1+\kappa_2>1$) at
$s=+\sqrt{3}$, where $S_{-,+}(s)$ exits the physical sheet and $S_{+,-}(s)$
enters from $P_+$, before coalescing with $S_{+,+}(s)$ at $s=2$. On the
reflectionless continuum phonon basis, we prove the \textbf{exact pointwise
identity}
\[
  \langle P_{D,2}\rangle_q + \langle P_{D,2}\rangle_{-q}
  \equiv \frac{q}{4\pi}\Big(\theta_D'(q;s,\eta) - \theta_D'(-q;s,\eta)\Big),
\]
resolving the open question in \texttt{audit\_charges\_action.pdf} by showing
directly from the quadratic boundary operator $P_{D,2}$ that
$\int_{-\infty}^\infty \dd q\,\langle P_{D,2}\rangle_q = \Delta P_D(q_D,\eta)$.
Similarly, Euler's theorem on the quadratic sewing Lagrangian $D_2 = a_t F_1 + b_t G_1 - E_{D,2}$
proves that $\langle E_{D,2}+J_{\mathrm{bulk},2}\rangle_q = \langle a_t F_1+b_t G_1\rangle_q$
and $-\big(\langle a_t F_1+b_t G_1\rangle_q+\langle a_t F_1+b_t G_1\rangle_{-q}\big) = \frac{\omega(q)}{4\pi}\big(\theta_D'(q)+\theta_D'(-q)\big)$.
\item \textbf{Boundary UV renormalisation and three sharp obstructions to physical quantum defect fusion (\sfus).}
With equal-mode-count cutoff $\Lambda$ in the asymptotic phase gauge
$\theta_D(\pm\infty)=0$, the Casimir energy $E_D^{\mathrm{mode}}(\Lambda)=E_D^{\mathrm{dens}}(\Lambda)-B_D(\Lambda)$
diverges linearly as $\frac{S_D}{\pi}(\sqrt{1+\Lambda^2}-1)+\Delta E_D^{\mathrm{ren}}+O(\Lambda^{-1})$
with $S_D = -\frac{\eta\sqrt{3}\sigma_{\mathrm{vac}}}{2} = 3\eta q_D$. Because
the bulk normal-ordering counterterm $\propto \int_{x\ne 0}\varphi_x\,\dd x$
vanishes on constant defect vacua, cancelling the defect UV divergence requires
either a distributional boundary tadpole counterterm $\propto (\varphi_R(0^+)-\varphi_L(0^-))$
or a boundary renormalisation of the defect parameters $(s,\rho_0,\tau_0,h_0)$.
Across KP triples we establish three obstructions to physical $O(\hb^0)$
triangle fusion:
\begin{itemize}
\item[(i)] \emph{Stability obstruction:} Every nondegenerate KP triple has
growing-mode counts $(0,1,2)$ generically (and $(0,1,1)$ at the branch-switch
locus $s_{\mathrm{eff}}=\sqrt{3}$), so \textbf{zero} single-soliton
absorption/emission edges connect two linearly stable defect vacua.
\item[(ii)] \emph{Canonical-charge holonomy no-go theorem:} While the mode-sum
charges $(\Delta E_1^{\mathrm{ren}},\Delta P_{\mathrm{mode}})$ have zero
directed sum around every KP triangle, the co-moving canonical
collective-coordinate charges $(\Delta E_{\mathrm{can}},\Delta P_{\mathrm{can}})$
have a \textbf{non-zero triangle holonomy}
$\sum_{i=1}^3\rho_i(\Delta E_{\mathrm{can}}(S_i),\Delta P_{\mathrm{can}}(S_i)) \ne (0,0)$.
Hence \emph{no} single-valued assignment of renormalised defect vacuum charges
$(E_D^{\mathrm{quant}}(q_i),P_D^{\mathrm{quant}}(q_i))$---regardless of boundary
counterterms---can balance $(\Delta E_{\mathrm{can}},\Delta P_{\mathrm{can}})$
at unshifted classical absorption rapidities across all three edges of a KP
triangle.
\item[(iii)] \emph{Absorbed-soliton momentum-renormalisation obstruction:}
Holding both soliton physical momenta $(P_3^{\mathrm{phys}},P_{S_i}^{\mathrm{phys}})$
fixed in $S(3,S_i)$ adds the second Taylor term
$\Delta X_{S_i}(3,S_i)\Delta P(S_i)$, whose directed sum around a KP triangle
equals $\sum_{i=1}^3\big(\frac{\Delta P(S_{i-1})}{k_{S_{i-1}}}-\frac{\Delta P(S_i)}{k_{S_i}}\big)\ln|z_t(3;q_i,\eta)| \ne 0$
for both $\Delta P_{\mathrm{mode}}$ and $\Delta P_{\mathrm{can}}$.
\end{itemize}
\end{enumerate}
\end{abstract}

\tableofcontents

% ==============================================================================
\section{Conventions, CZ dictionary, and scope}
\label{sec:setup}
% ==============================================================================

We work throughout in the fixed project conventions (\texttt{AGENTS.md},
\texttt{README.md}) for the \textbf{good} Boussinesq equation \La:
\begin{equation}\label{eq:good_bq}
  u_{tt} = u_{xx} - u_{xxxx} - 3(u^2)_{xx}, \qquad
  u = \varphi_x = 2(\ln\tau)_{xx},
\end{equation}
with bulk Lagrangian density on the left ($x<0$, $\varphi_L$) and right ($x>0$,
$\varphi_R$) half-lines and canonical symplectic form
\begin{equation}\label{eq:bulk_action}
  \mathcal{L}(\varphi) = \frac{1}{2}\varphi_t^2 - \frac{1}{2}\varphi_x^2
  + \varphi_x^3 - \frac{1}{2}\varphi_{xx}^2, \qquad
  \Omega_{\mathrm{bulk}} = \int_{-\infty}^0 \dd\pi_L\wedge\dd\varphi_L
  + \int_0^\infty \dd\pi_R\wedge\dd\varphi_R, \quad \pi_{L,R} = \varphi_{L,R,t}.
\end{equation}
Planck's constant $\hb$ refers to this normalisation.

A one-soliton with amplitude parameter $\alpha \in (0,\pi/2)$ and propagation
direction $\epsilon \in \{+1,-1\}$ has \La/\Lc:
\begin{align}\label{eq:sol_data}
  k &= \sin\alpha, \quad v = \epsilon\cos\alpha, \quad \ell = -\epsilon\cos 3\alpha, \quad
  M = 2k, \nonumber\\
  \Pi &= -Mv, \quad P = \frac{2}{3}\epsilon k^3\cos\alpha, \quad
  E = \frac{2}{15}k^3(5-4k^2),
\end{align}
and KP roots $p_a > p_b$ satisfying $p_a - p_b = k$, $p_a + p_b = v/\sqrt{3}$,
and $4p^3 - p = -\frac{\epsilon}{3\sqrt{3}}\cos 3\alpha = \frac{\ell}{3\sqrt{3}}$
\Lc\ (\texttt{defect\_semiclassics.pdf}).

\paragraph{Corrigan--Zambon dictionary.}
Corrigan and Zambon (CZ) \cite{CZ} (with $c=\lambda=\kappa=1$) write $u_{\mathrm{CZ}}$
for $x<0$ and $v_{\mathrm{CZ}}$ for $x>0$ with bulk Lagrangian
$\mathcal{L}_{\mathrm{CZ}} = \frac{1}{2}u_{\mathrm{CZ},t}^2 - \frac{1}{2}u_{\mathrm{CZ},x}^2 - \frac{1}{3}u_{\mathrm{CZ},x}^3 - \frac{1}{2}u_{\mathrm{CZ},xx}^2$.
The explicit dictionary to project variables is \Lc\ (\sclass):
\begin{equation}\label{eq:cz_dict}
  u_{\mathrm{CZ}} = -3\varphi_L, \qquad v_{\mathrm{CZ}} = -3\varphi_R, \qquad
  \mathcal{L}_{\mathrm{CZ}} = 9\,\mathcal{L}, \qquad
  D_{\mathrm{proj}} = \frac{1}{9}D_{\mathrm{CZ}}, \qquad
  (E_{\mathrm{CZ}},P_{\mathrm{CZ}}) = 9\,(E,P).
\end{equation}
We do not rely on any results from Phase~2i (\texttt{toda\_wh\_and\_kp\_synthesis.pdf}),
which is under separate audit.

% ==============================================================================
\section{Classical localised defect charges and exact conservation laws}
\label{sec:classical}
% ==============================================================================

\subsection{Sewing conditions and localised charges in project conventions}

In CZ variables \cite{CZ} \La, the type-I defect at $x=0$ has orientation
$\eta \in \{+1,-1\}$, $\alpha_{\mathrm{CZ}} = \eta/\sqrt{3}$,
$\beta_{\mathrm{CZ}} = 2\alpha_{\mathrm{CZ}} = 2\eta/\sqrt{3}$,
$\mu = \frac{2}{3\beta_{\mathrm{CZ}}} = \eta/\sqrt{3}$, and bare parameter
$s = \sigma_0 \in \mathbb{R}$ (with KP vertex $q_D = -\frac{s}{2\sqrt{3}}$).
In terms of the potential jump at $x=0$,
\begin{equation}\label{eq:w_phi_def}
  w_\varphi := \varphi_L(0^-,t) - \varphi_R(0^+,t) = -M_D, \qquad
  w_{\mathrm{CZ}} = u_{\mathrm{CZ}}(0^-,t) - v_{\mathrm{CZ}}(0^+,t) = -3w_\varphi = 3M_D,
\end{equation}
the dynamical defect field $\sigma$ and auxiliary polynomials $\rho(\sigma)$,
$\tau(\sigma)$, and $h(\sigma)$ are \La/\Lc\ (\sclass):
\begin{align}
  \sigma(w_\varphi) &= s + \mu w_{\mathrm{CZ}} = s - \sqrt{3}\,\eta\,w_\varphi, \qquad
  \rho_0 = -\frac{\sqrt{3}\,\eta}{2}\,s\!\left(1 - \frac{s^2}{3}\right) = \frac{\eta}{\sqrt{3}}\,\ell_D, \label{eq:sigma_rho0}\\
  \rho(\sigma) &= \rho_0 + \frac{\sqrt{3}\,\eta}{2}\!\left(\sigma - \frac{\sigma^3}{3}\right)
  = \frac{\sqrt{3}\,\eta}{6}\,(\sigma - s)\big(3 - (\sigma^2 + s\sigma + s^2)\big), \label{eq:rho_def}\\
  \tau(\sigma) &= -\int_s^\sigma \frac{\beta_{\mathrm{CZ}}}{2\mu}\,x\,\rho(x)\,\dd x
  = -\int_s^\sigma x\,\rho(x)\,\dd x, \label{eq:tau_def}\\
  h(\sigma) &= -\int_s^\sigma \frac{\beta_{\mathrm{CZ}}-\alpha_{\mathrm{CZ}}}{\mu}\,\rho(x)\,\dd x
  = -\int_s^\sigma \rho(x)\,\dd x. \label{eq:h_def}
\end{align}
In the shift-invariant gauge $f = g = -\frac{\sigma^2}{4\mu} = -\frac{\sqrt{3}\,\eta}{4}\sigma^2$
(which satisfies $f_{v_{\mathrm{CZ}}} - g_{u_{\mathrm{CZ}}} = \sigma$ and preserves
$\varphi \mapsto \varphi + c$ \Lc, \texttt{defect\_semiclassics.pdf}), the
project-normalised defect Lagrangian at $x=0$ is
\begin{equation}\label{eq:D_proj}
  D_{\mathrm{proj}} = \varphi_{L,t} F_{\mathrm{proj}} + \varphi_{R,t} G_{\mathrm{proj}} - E_D,
\end{equation}
where \Lc\ (\sclass):
\begin{align}
  F_{\mathrm{proj}} &= \frac{\eta}{\sqrt{3}}\big(\varphi_{L,x} + 2\varphi_{R,x}\big) + \frac{\sqrt{3}\,\eta}{12}\,\sigma^2, \qquad
  G_{\mathrm{proj}} = \frac{\eta}{\sqrt{3}}\big(2\varphi_{L,x} + \varphi_{R,x}\big) + \frac{\sqrt{3}\,\eta}{12}\,\sigma^2, \label{eq:FG_proj}\\
  E_D &= -\frac{H_{\mathrm{CZ}}}{9}
  = -\frac{\eta}{\sqrt{3}}\,\sigma\big(\varphi_{L,x}^2 + \varphi_{L,x}\varphi_{R,x} + \varphi_{R,x}^2\big)
    + \frac{1}{3}\,\rho(\sigma)\big(\varphi_{L,x} + \varphi_{R,x}\big) - \frac{1}{9}\,\tau(\sigma). \label{eq:ED_proj}
\end{align}

\begin{proposition}[Classical localised defect charges and off-shell conservation, \Lc, \sclass]\label{prop:classical_charges}
Varying $S = \int_{x<0}\mathcal{L}(\varphi_L)\,\dd x\,\dd t + \int_{x>0}\mathcal{L}(\varphi_R)\,\dd x\,\dd t + \int D_{\mathrm{proj}}\,\dd t$
yields the four sewing conditions at $x=0$:
\begin{align}
  \varphi_{L,xx} &= \frac{\partial D_{\mathrm{proj}}}{\partial \varphi_{L,x}}, &
  \varphi_{L,xxx} &= \varphi_{L,x} - 3\varphi_{L,x}^2 - \frac{\partial D_{\mathrm{proj}}}{\partial \varphi_L} + \frac{\dd F_{\mathrm{proj}}}{\dd t}, \label{eq:sew_L}\\
  \varphi_{R,xx} &= -\frac{\partial D_{\mathrm{proj}}}{\partial \varphi_{R,x}}, &
  \varphi_{R,xxx} &= \varphi_{R,x} - 3\varphi_{R,x}^2 + \frac{\partial D_{\mathrm{proj}}}{\partial \varphi_R} - \frac{\dd G_{\mathrm{proj}}}{\dd t}. \label{eq:sew_R}
\end{align}
For any smooth solution of the bulk equation \eqref{eq:good_bq} on $x \ne 0$
satisfying \eqref{eq:sew_L}--\eqref{eq:sew_R} at $x=0$, the four total charges
\begin{equation}\label{eq:tot_charges}
  E_{\mathrm{tot}} = E_{\mathrm{bulk}} + E_D, \quad
  P_{\mathrm{tot}} = P_{\mathrm{bulk}} + P_D, \quad
  \Pi_{\mathrm{tot}} = \Pi_{\mathrm{bulk}} + \Pi_D, \quad
  M_{\mathrm{tot}} = M_{\mathrm{bulk}} + M_D,
\end{equation}
are conserved identically ($\frac{\dd}{\dd t}(E_{\mathrm{tot}},P_{\mathrm{tot}},\Pi_{\mathrm{tot}},M_{\mathrm{tot}}) = (0,0,0,0)$),
where $E_D$ is given by \eqref{eq:ED_proj} and the localised defect momentum
$P_D$, shift charge $\Pi_D$, and mass jump $M_D$ in project normalisation are:
\begin{align}
  P_D &= \frac{P_{D,\mathrm{CZ}}}{9}
  = -\frac{\eta}{2\sqrt{3}}\big(\varphi_{L,x}^2 + 4\varphi_{L,x}\varphi_{R,x} + \varphi_{R,x}^2\big) + \frac{1}{9}\,h(\sigma), \label{eq:PD_proj}\\
  \Pi_D &= F_{\mathrm{proj}} + G_{\mathrm{proj}}
  = \sqrt{3}\,\eta\big(\varphi_{L,x} + \varphi_{R,x}\big) + \frac{\sqrt{3}\,\eta}{6}\,\sigma^2, \label{eq:PiD_proj}\\
  M_D &= \varphi_R(0^+,t) - \varphi_L(0^-,t) = -w_\varphi = \frac{\sigma - s}{\sqrt{3}\,\eta}. \label{eq:MD_proj}
\end{align}
\end{proposition}

\begin{remark}[Avoiding double-subtraction in $P_D$, \Lc]
In CZ \cite[(3.11)]{CZ}, the modified momentum is introduced in the intermediate
form $P_{\mathrm{bulk},\mathrm{CZ}} - u_{\mathrm{CZ},x}F_{\mathrm{CZ}} - v_{\mathrm{CZ},x}G_{\mathrm{CZ}} + \tilde{P}_D$,
and in their final formula \cite[(3.49)]{CZ} the derivative terms
$-u_{\mathrm{CZ},x}F_{\mathrm{CZ}} - v_{\mathrm{CZ},x}G_{\mathrm{CZ}}$ have
already been combined with $\tilde{P}_D$ to produce
$P_{D,\mathrm{CZ}} = -\frac{\alpha_{\mathrm{CZ}}}{2}(u_{\mathrm{CZ},x}^2+v_{\mathrm{CZ},x}^2) - \beta_{\mathrm{CZ}}u_{\mathrm{CZ},x}v_{\mathrm{CZ},x} + h(\sigma)$,
which is free of time derivatives. Subtracting $u_{\mathrm{CZ},x}F_{\mathrm{CZ}} + v_{\mathrm{CZ},x}G_{\mathrm{CZ}}$
a second time leaves a non-zero time-derivative remainder (recorded in
Section~\ref{sec:failures}, Failure~1).
\end{remark}

\subsection{Exact conservation through soliton absorption \texorpdfstring{$D(q_i) + S_i \to D(q_{i+1})$}{D(q\_i) + S\_i -> D(q\_{i+1})}}

On a static defect vacuum ($\varphi_{L,x} = \varphi_{R,x} = 0$), the sewing
condition \eqref{eq:sew_L} requires $\rho(\sigma) = 0$, whose three roots
$\sigma_i = -2\sqrt{3}\,q_i$ correspond to the three roots $\{q_1,q_2,q_3\}$ of
the KP cubic $q_i^2 + q_i q_j + q_j^2 = \frac{1}{4}$ containing $q_D = -\frac{s}{2\sqrt{3}}$
\Lc\ (\texttt{defect\_semiclassics.pdf}). Evaluating \eqref{eq:ED_proj} and
\eqref{eq:PD_proj}--\eqref{eq:MD_proj} on $\varphi_{L,x}=\varphi_{R,x}=0$ gives
the vacuum charges of vertex $q_i$ \Lc\ (\sclass):
\begin{equation}\label{eq:vac_charges}
  E_D^{\mathrm{vac}}(q_i) = -\frac{1}{9}\tau(\sigma_i), \;
  P_D^{\mathrm{vac}}(q_i) = \frac{1}{9}h(\sigma_i), \;
  \Pi_D^{\mathrm{vac}}(q_i) = \frac{\sqrt{3}\,\eta}{6}\sigma_i^2 = 2\sqrt{3}\,\eta\,q_i^2, \;
  M_D^{\mathrm{vac}}(q_i) = 2\eta(q_D - q_i).
\end{equation}

\begin{proposition}[Symbolic and 50-digit time-dependent conservation in absorption, \Lc, \sclass]\label{prop:classical_absorption}
Let $\{q_1,q_2,q_3\}$ be a KP triple ($q_i^2 + q_i q_{i+1} + q_{i+1}^2 = \frac{1}{4}$),
$S_i = \{q_i,q_{i+1}\}$ the soliton with KP roots $p_a = \max(q_i,q_{i+1})$ and
$p_b = \min(q_i,q_{i+1})$, and $\rho_i := \eta\sgn(q_i - q_{i+1}) \in \{+1,-1\}$.
\begin{enumerate}
\item Modulo the KP triple polynomial $q_i^2 + q_i q_{i+1} + q_{i+1}^2 - \frac{1}{4} = 0$,
the vacuum jumps of all four defect charges \eqref{eq:vac_charges} equal $\rho_i$
times the classical soliton charges \eqref{eq:sol_data} identically:
\begin{equation}\label{eq:vac_jump_all4}
  \boxed{\;
  \big(E_D^{\mathrm{vac}},P_D^{\mathrm{vac}},\Pi_D^{\mathrm{vac}},M_D^{\mathrm{vac}}\big)\big|_{q_i}^{q_{i+1}}
  = \rho_i\,\big(E(S_i),\,P(S_i),\,\Pi(S_i),\,M(S_i)\big).
  \;}
\end{equation}
\item On the exact time-dependent CZ one-soliton solutions---both in regular
transmission ($z>0$) and in exact absorption $D(q_1)+S_1 \to D(q_2)$ ($z=0$ on
$x>0$)---direct spatial quadrature of $(E_{\mathrm{bulk}}(t),P_{\mathrm{bulk}}(t),\Pi_{\mathrm{bulk}}(t),M_{\mathrm{bulk}}(t))$
plus the localised defect charges $(E_D(t),P_D(t),\Pi_D(t),M_D(t))$ at $x=0$ is
independent of $t \in [-8,+8]$ to $9.35\times 10^{-51}$ (transmission) and
$5.35\times 10^{-51}$ (absorption).
\end{enumerate}
\end{proposition}

% ==============================================================================
\section{One-loop defect charges, spectral branches, and boundary renormalisation}
\label{sec:oneloop}
% ==============================================================================

\subsection{Corrected sewing determinant and the admissible branch switch at \texorpdfstring{$s_{\mathrm{eff}} = \sqrt{3}$}{s\_eff = sqrt(3)}}

Linearising the sewing conditions \eqref{eq:sew_L}--\eqref{eq:sew_R} around a
static defect vacuum $\sigma_{\mathrm{vac}}$ with effective parameter
$s_{\mathrm{eff}} := \eta\sigma_{\mathrm{vac}} = -2\sqrt{3}\,\eta\,q_{\mathrm{vac}}$
and seeking localised growing modes of frequency $\omega = \ii g$ ($g>0$) and
spatial decay exponents $\kappa_{1,2} > 0$ ($\kappa_1^2+\kappa_2^2=1$,
$\kappa_1\kappa_2=g$, $S := \kappa_1+\kappa_2 = \sqrt{1+2g} > 1$,
$d := \kappa_1-\kappa_2 = \sqrt{2-S^2}$) yields a $4\times 4$ sewing matrix
$M_{4\times 4}(g;\sigma_{\mathrm{vac}},\eta)$. Following the independent audit
\texttt{audit\_frames\_defects.pdf}, we verify its exact off-shell polynomial
factorisation and determine the complete admissible branch structure \Lc\
(\sren).

\begin{proposition}[Corrected sewing determinant and admissible branch switch, \Lc, \sren]\label{prop:det_branches}
\begin{enumerate}
\item Off-shell for all $(S,s_{\mathrm{eff}})$, the reduced $4\times 4$
linearised sewing determinant is
\begin{equation}\label{eq:det_corrected}
  \frac{\det M_{4\times 4}(g;\sigma_{\mathrm{vac}},\eta)}{(\kappa_1-\kappa_2)^2}
  = -\frac{2}{3}\,S^2(S^2-1)\,P_-(S;s_{\mathrm{eff}})\,P_+(S;s_{\mathrm{eff}}),
\end{equation}
where $P_\mp(S;s_{\mathrm{eff}}) = S^2 - \frac{\sqrt{3}\mp 1}{2}s_{\mathrm{eff}}S + \frac{s_{\mathrm{eff}}^2-2\mp\sqrt{3}}{2}$.
(The coefficient $-\frac{1}{24}$ printed in Phase~2g is sixteen times too small
because column replacement $M_{\cdot,1}\pm M_{\cdot,2}$ multiplies the
determinant by $4$.)
\item At the continuum threshold $S = 1$ ($g = 0$), the two quadratic factors
factorise exactly as
\begin{equation}\label{eq:P_at_S1}
  P_-(1;s_{\mathrm{eff}}) = \frac{1}{2}(s_{\mathrm{eff}}+1)(s_{\mathrm{eff}}-\sqrt{3}), \qquad
  P_+(1;s_{\mathrm{eff}}) = \frac{1}{2}(s_{\mathrm{eff}}-1)(s_{\mathrm{eff}}-\sqrt{3}).
\end{equation}
\item Retaining only admissible decaying half-line roots ($S > 1$ for
$s_{\mathrm{eff}} \in (-2,2)$):
\begin{itemize}
\item For $s_{\mathrm{eff}} \le -1$ ($\eta q_{\mathrm{vac}} \ge \frac{1}{2\sqrt{3}}$):
\textbf{zero} growing modes exist; these are the only linearly stable defect
vacua, and the only vacua admissible as external quantum states.
\item For $-1 < s_{\mathrm{eff}} \le 1$: exactly \textbf{one} growing mode exists,
given by $S_{-,+}(s_{\mathrm{eff}}) > 1$ from $P_- = 0$, with growth rate
$g_-(s_{\mathrm{eff}}) = \frac{s_{\mathrm{eff}}\sqrt{4-s_{\mathrm{eff}}^2}+\sqrt{3}(2-s_{\mathrm{eff}}^2)}{8} > 0$.
\item For $1 < s_{\mathrm{eff}} < 2$: a second growing mode $S_{+,+}(s_{\mathrm{eff}}) > 1$
emerges from $P_+ = 0$ at $s_{\mathrm{eff}} = 1$ with growth rate
$g_{+,+}(s_{\mathrm{eff}}) = \frac{s_{\mathrm{eff}}\sqrt{4-s_{\mathrm{eff}}^2}-\sqrt{3}(2-s_{\mathrm{eff}}^2)}{8} > 0$.
At $s_{\mathrm{eff}} = \sqrt{3}$, both $P_-(1;\sqrt{3})=0$ and $P_+(1;\sqrt{3})=0$:
the root $S_{-,+}(s_{\mathrm{eff}})$ of $P_-$ decreases through $S=1$ ($g=0$) and
leaves the admissible sheet $S>1$, while the second root
$S_{+,-}(s_{\mathrm{eff}}) = \frac{(\sqrt{3}+1)s_{\mathrm{eff}}-(\sqrt{3}-1)\sqrt{4-s_{\mathrm{eff}}^2}}{4}$
of $P_+$ simultaneously increases through $S=1$ and enters the admissible sheet
for $\sqrt{3} < s_{\mathrm{eff}} < 2$. Thus on $(1,2)$ the two non-negative
admissible growth rates are (verified to $55$ digits)
\begin{equation}\label{eq:g_admissible}
  g_1(s_{\mathrm{eff}}) = \frac{s_{\mathrm{eff}}\sqrt{4-s_{\mathrm{eff}}^2} - \sqrt{3}(2-s_{\mathrm{eff}}^2)}{8} > 0, \quad
  g_2(s_{\mathrm{eff}}) = \frac{\big|s_{\mathrm{eff}}\sqrt{4-s_{\mathrm{eff}}^2} + \sqrt{3}(2-s_{\mathrm{eff}}^2)\big|}{8} \ge 0,
\end{equation}
which both belong to $P_+(S;s_{\mathrm{eff}})=0$ for $\sqrt{3}<s_{\mathrm{eff}}\le 2$
and coalesce at $s_{\mathrm{eff}}=2$ to $g_1(2)=g_2(2)=\frac{\sqrt{3}}{4}$.
\end{itemize}
\end{enumerate}
\end{proposition}

\subsection{Direct operator derivation of \texorpdfstring{$\Delta P_D$ and $\Delta E_D$ from $P_{D,2}$ and $D_2$}{Delta P\_D and Delta E\_D from P\_{D,2} and D\_2}}

Expanding the project-normalised defect Lagrangian \eqref{eq:D_proj} and
localised momentum \eqref{eq:PD_proj} to quadratic order in potential
fluctuations $(a,b) = (\delta\varphi_L,\delta\varphi_R)$ about a constant vacuum
$\sigma_0$ gives \Lc\ (\texttt{audit\_charges\_action.pdf}, \sren):
\begin{align}
  D_2 &= a_t F_1 + b_t G_1 - E_{D,2}, \qquad
  F_1 = \alpha a_x + \beta b_x - \frac{\sigma_0}{2}(a-b), \qquad
  G_1 = \beta a_x + \alpha b_x - \frac{\sigma_0}{2}(a-b), \label{eq:D2_F1G1}\\
  E_{D,2} &= -H_2 = -\sigma_0\!\left[\frac{\beta}{2}(a_x^2+b_x^2) + \alpha a_x b_x\right]
  - \frac{1-\sigma_0^2}{2}(a-b)(a_x+b_x) + \frac{\beta\sigma_0(1-\sigma_0^2)}{8}(a-b)^2, \label{eq:ED2}\\
  P_{D,2} &= -\frac{\alpha}{2}(a_x^2+b_x^2) - \beta a_x b_x - \frac{\alpha}{4}(1-\sigma_0^2)(a-b)^2, \label{eq:PD2}
\end{align}
where $\alpha = \eta/\sqrt{3}$ and $\beta = 2\alpha$. In
\texttt{audit\_charges\_action.pdf} (Section~4), it was noted as an open issue
\Ld\ that staticity alone does not prove why the expectation value of the
conserved momentum including $P_{D,2}$ equals the phase-shift integral
$\int_{-\infty}^\infty \frac{\dd q}{4\pi}\,q\,\theta_D'(q)$. We now supply this
derivation \Lc.

On any constant defect vacuum, the continuum phonon eigenmodes are reflectionless
plane waves all the way to $x=0^\mp$ (with zero evanescent tails \Lc, Claim~2.1
of \texttt{oneloop\_\allowbreak defect\_\allowbreak phase.pdf}). With canonical bulk field normalisation
($[\varphi(x),\pi(y)] = \ii\hb\delta(x-y)$, zero-point weight $\frac{1}{2}\hb$ per
mode), a mode of wavenumber $q \in \mathbb{R}\setminus\{0\}$ and frequency
$\omega(q) = |q|\sqrt{1+q^2}$ has
\begin{align}\label{eq:canon_phonon}
  a_q(x,t) &= \frac{1}{\sqrt{\pi\omega(q)}}\,\Rea\!\big(\ee^{\ii(qx-\omega(q)t)}\big) \quad (x<0), \nonumber\\
  b_q(x,t) &= \frac{1}{\sqrt{\pi\omega(q)}}\,\Rea\!\big(z_{\mathrm{ph}}(q;q_D,\eta)\,\ee^{\ii(qx-\omega(q)t)}\big) \quad (x>0),
\end{align}
so that $\langle a^2\rangle_q = \langle b^2\rangle_q = \frac{1}{4\pi\omega(q)}$,
$\langle a_x^2\rangle_q = \langle b_x^2\rangle_q = \frac{q^2}{4\pi\omega(q)}$,
$\langle a_x b_x\rangle_q = \frac{q^2\Rea z_{\mathrm{ph}}(q)}{4\pi\omega(q)}$,
$\langle (a-b)^2\rangle_q = \frac{2(1-\Rea z_{\mathrm{ph}}(q))}{4\pi\omega(q)}$,
and $\langle (a-b)(a_x+b_x)\rangle_q = -\frac{2q\Ima z_{\mathrm{ph}}(q)}{4\pi\omega(q)}$
at $x=0$. Because $|z_{\mathrm{ph}}(q)| = 1$ everywhere on $x>0$, the bulk
momentum density $-\varphi_t\varphi_x$ on mode $q$ is spatially constant and
identical to its free-vacuum value, so its vacuum-subtracted bulk integral and
its $\pm q$ symmetric bulk sum vanish identically. All vacuum-subtracted
momentum of the mode pair $(+q,-q)$ is therefore localised in $P_{D,2}$ at $x=0$.

\begin{proposition}[Exact pointwise identity between $\langle P_{D,2}\rangle_q$ and $\frac{q}{4\pi}\theta_D'(q)$, \Lc, \sren]\label{prop:PD2_identity}
For every $q > 0$, every $\sigma_0 \in \mathbb{R}$, and $\eta \in \{+1,-1\}$:
\begin{enumerate}
\item As $q \to \infty$, $z_{\mathrm{ph}}(q;q_D,\eta) \to \ee^{-\ii\eta\sgn(q)\frac{2\pi}{3}}$
 has $\Rea z_{\mathrm{ph}}(q) \to -\frac{1}{2}$. Because $\beta = 2\alpha$, the
leading $O(q^0)$ derivative combination $-\frac{\alpha}{2}(a_x^2+b_x^2) - \beta a_x b_x \propto 1 + 2\Rea z_{\mathrm{ph}}(q) \to 0$
cancels identically at $O(1)$, leaving $\langle P_{D,2}\rangle_q = -\frac{\sqrt{3}\eta\sigma_0}{8\pi q} + \frac{\sqrt{3}\eta(\sigma_0^2-2)}{16\pi q^2} + O(q^{-3})$.
\item Symmetrising over $+q$ and $-q$ cancels the odd $O(q^{-1})$ tail and gives
the \textbf{exact pointwise algebraic identity}
\begin{equation}\label{eq:PD2_pointwise}
  \boxed{\;
  \begin{aligned}
  \langle P_{D,2}\rangle_q + \langle P_{D,2}\rangle_{-q}
  &\equiv \frac{q}{4\pi}\Big(\theta_D'(q;s,\eta) - \theta_D'(-q;s,\eta)\Big) \\
  &= \frac{\sqrt{3}\,\eta\,q\big(2(\sigma_0^2-2)q^2 + \sigma_0^2 - 1\big)}
         {\pi\sqrt{1+q^2}\big(16q^4 + 4(\sigma_0^2+2)q^2 + (\sigma_0^2-1)^2\big)}.
  \end{aligned}
  \;}
\end{equation}
Consequently, the vacuum expectation value of the localised quadratic defect
momentum operator is UV finite and equals $\Delta P_D(q_D,\eta)$ identically:
\begin{equation}\label{eq:PD2_integral}
  \langle P_{D,2}\rangle_{\mathrm{vac}} := \int_{-\infty}^\infty \dd q\,\langle P_{D,2}\rangle_q
  = \int_{-\infty}^\infty \frac{\dd q}{4\pi}\,q\,\theta_D'(q;s,\eta) = \Delta P_D(q_D,\eta).
\end{equation}
\item For the quadratic energy, integrating the bulk potential energy
$\frac{1}{2}\int_{x\ne 0}(\varphi_x^2+\varphi_{xx}^2)\,\dd x$ by parts at $x=0^\mp$
produces the boundary jump $J_{\mathrm{bulk},2} := \frac{1}{2}[a(a_x-a_{xxx})+a_x a_{xx}]_{0^-} - \frac{1}{2}[b(b_x-b_{xxx})+b_x b_{xx}]_{0^+}$.
By Euler's homogeneous function theorem on $D_2 = a_t F_1 + b_t G_1 - E_{D,2}$
and the linearised sewing equations, every harmonic sewing mode satisfies
$\langle E_{D,2} + J_{\mathrm{bulk},2}\rangle_q = \langle a_t F_1 + b_t G_1\rangle_q$,
and the symmetric sum over $\pm q$ obeys the \textbf{exact pointwise identity}
\begin{equation}\label{eq:ED2_pointwise}
  \boxed{\;
  -\Big(\langle a_t F_1 + b_t G_1\rangle_q + \langle a_t F_1 + b_t G_1\rangle_{-q}\Big)
  \equiv \frac{\omega(q)}{4\pi}\Big(\theta_D'(q;s,\eta) + \theta_D'(-q;s,\eta)\Big).
  \;}
\end{equation}
\end{enumerate}
\end{proposition}

\subsection{Cutoff Casimir energy, boundary UV divergence, and counterterm analysis}

On any linearly stable defect vacuum ($s_{\mathrm{eff}} = \eta\sigma_{\mathrm{vac}} < -1$),
there are no bound or growing discrete modes. In the equal-mode-count Gaussian
quantisation prescription \cite[App.~6]{FK} (audited for the bulk soliton in
\texttt{audit\_charges\_action.pdf}), integrating the mode-shift sum by parts on
$[-\Lambda,0^-] \cup [0^+,\Lambda]$ requires the \textbf{asymptotic scattering
branch} $\theta_D^{\mathrm{asympt}}(q;s,\eta) := \sum_{s'=1}^2 c_{D,s'}\Theta(q;\psi_{s'})$,
which satisfies $\theta_D^{\mathrm{asympt}}(\pm\infty;s,\eta) = 0$ (whereas the
zero-endpoint branch $\theta_D^{\mathrm{phys}}(q)$ with $\theta_D^{\mathrm{phys}}(0^\pm)=0$
approaches $\pm\pi\eta$ at $q\to\pm\infty$ for $s_{\mathrm{eff}}<-1$; its nonzero
jump at $q=0^\pm$ is multiplied by $\omega(0)=0$ in the energy integration by
parts and therefore drops out \Lc, Section~\ref{sec:failures}, Failure~3).

\begin{proposition}[Regulated defect Casimir energy and boundary counterterm structure, \Lc, \sren]\label{prop:cutoff_renorm}
Let $s_{\mathrm{eff}} = \eta\sigma_{\mathrm{vac}} < -1$,
$S_D := \sum_{s'=1}^2 c_{D,s'}\sin\psi_{s'} = -\frac{\eta\sqrt{3}\sigma_{\mathrm{vac}}}{2} = 3\eta q_{\mathrm{vac}} > \frac{\sqrt{3}}{2}$,
and $T := \sqrt{1+\Lambda^2}$. Then:
\begin{enumerate}
\item As $\Lambda \to \infty$, the phase-density integral $E_D^{\mathrm{dens}}(\Lambda)$
and cutoff endpoint $B_D(\Lambda)$ in the asymptotic branch $\theta_D^{\mathrm{asympt}}(\pm\infty)=0$
satisfy (verified to $<10^{-20}$ at $\Lambda=10^{22}$)
\begin{align}
  E_D^{\mathrm{dens}}(\Lambda) &:= \int_{-\Lambda}^\Lambda \frac{\dd q}{4\pi}\,\omega(q)\theta_D'(q)
  = \frac{S_D}{2\pi}(T - 1) + \left(\Delta E_D^{\mathrm{ren}} - \frac{S_D}{2\pi}\right) + O(\Lambda^{-1}), \label{eq:ED_dens}\\
  B_D(\Lambda) &:= \frac{\Lambda\sqrt{1+\Lambda^2}}{4\pi}\Big[\theta_D^{\mathrm{asympt}}(\Lambda) - \theta_D^{\mathrm{asympt}}(-\Lambda)\Big]
  = -\frac{S_D}{2\pi}\,T + O(\Lambda^{-1}), \label{eq:BD_end}\\
  E_D^{\mathrm{mode}}(\Lambda) &:= E_D^{\mathrm{dens}}(\Lambda) - B_D(\Lambda)
  = \frac{S_D}{\pi}(T - 1) + \Delta E_D^{\mathrm{ren}} + O(\Lambda^{-1}), \label{eq:ED_mode}
\end{align}
where $\frac{T-1}{2\pi} = \int_{-\Lambda}^\Lambda \frac{\dd q}{2\pi}\frac{|q|}{2\sqrt{1+q^2}} = \langle\varphi_x^2\rangle_{\mathrm{vac}}(\Lambda)$
is the exact one-loop bulk tadpole integral and $\Delta E_D^{\mathrm{ren}}$ is
the finite 2-root dipole formula of Phase~2g. The momentum endpoint
$\frac{\Lambda}{4\pi}[\theta_D^{\mathrm{asympt}}(\Lambda)+\theta_D^{\mathrm{asympt}}(-\Lambda)]$
vanishes as $O(\Lambda^{-1})$.
\item \textbf{Need for a boundary counterterm:} Because the open-half-line bulk
normal-ordering counterterm $\delta E_{\mathrm{bulk}}^{\mathrm{ct}} = \frac{3(T-1)}{2\pi}\big(\int_{-\infty}^0 \varphi_{L,x}\,\dd x + \int_0^\infty \varphi_{R,x}\,\dd x\big)$
vanishes identically on every spatially constant defect vacuum ($\varphi_{L,x}^{(0)}=\varphi_{R,x}^{(0)}=0$),
cancelling the UV divergence $-\frac{S_D}{\pi}(T-1) = \frac{\eta\sqrt{3}\sigma_{\mathrm{vac}}}{2\pi}(T-1)$
of $E_D^{\mathrm{mode}}(\Lambda)$ requires an explicit localised boundary
counterterm at $x=0$:
\begin{equation}\label{eq:delta_ED_ct}
  \delta E_D^{\mathrm{ct}}(\sigma_{\mathrm{vac}};\Lambda)
  = -\frac{S_D}{\pi}(T - 1) + \delta E_{D,\mathrm{fin}}(\sigma_{\mathrm{vac}})
  = \frac{\eta\sqrt{3}(T-1)}{2\pi}\big(s + \sqrt{3}\,\eta\,M_D\big) + \delta E_{D,\mathrm{fin}}(\sigma_{\mathrm{vac}}).
\end{equation}
Notice that the $M_D$-dependent part of \eqref{eq:delta_ED_ct},
$+\frac{3(T-1)}{2\pi}M_D = \frac{3(T-1)}{2\pi}(\varphi_R(0^+)-\varphi_L(0^-))$,
is precisely the boundary jump required to complete the open-half-line bulk
integral $\int_{x\ne 0}\varphi_x\,\dd x = (\varphi_R(+\infty)-\varphi_L(-\infty)) - M_D$
into the topological total-derivative counterterm $\frac{3(T-1)}{2\pi}(\varphi_R(+\infty)-\varphi_L(-\infty))$!
\item \textbf{Integrability constraint on finite boundary counterterms:}
Adding a shift-invariant zero-derivative potential $\delta\tau_{\mathrm{fin}}(w_\varphi)$
to $H_{\mathrm{CZ}}$ without modifying the derivative sewing terms changes the
momentum flux balance on the sewing manifold by
$-(u_{\mathrm{CZ},x}-v_{\mathrm{CZ},x})\delta\tau_{\mathrm{fin}}'(w_{\mathrm{CZ}})$,
which cannot equal a total time derivative $\frac{\dd}{\dd t}\delta h_{\mathrm{fin}}(w_{\mathrm{CZ}}) = (u_{\mathrm{CZ},t}-v_{\mathrm{CZ},t})\delta h_{\mathrm{fin}}'(w_{\mathrm{CZ}})$
for independent boundary jets unless $\delta\tau_{\mathrm{fin}}' = \delta h_{\mathrm{fin}}' = 0$.
Thus an off-shell integrable one-loop boundary renormalisation must take the form
of a one-loop shift of the integrable CZ defect parameters
\begin{equation}\label{eq:param_renorm}
  (s,\,\rho_0,\,\tau_0,\,h_0) \;\longmapsto\;
  (s + \hb\,\delta s_1,\;\rho_0 + \hb\,\delta\rho_{0,1},\;\tau_0 + \hb\,\delta\tau_{0,1},\;h_0 + \hb\,\delta h_{0,1}).
\end{equation}
Here $(\delta\tau_{0,1},\delta h_{0,1})$ are vacuum-independent additive constants
that cancel in vacuum jumps; keeping $w_\varphi = 0$ ($\sigma = s$) as a static
vacuum requires $\rho(s)=0$, which ties $\delta\rho_{0,1}$ to $\delta s_1$ and
leaves \textbf{one free finite parameter renormalisation $\delta s_1$} (or two
parameters $(\delta s_1,\delta\rho_{0,1})$ if the vacuum jump $w_\varphi^{(0)} = O(\hb)$
is allowed to shift at one loop).
\end{enumerate}
\end{proposition}

% ==============================================================================
\section{Quantum absorption and fusion test at \texorpdfstring{$O(\hb^0)$}{O(hbar\^0)}}
\label{sec:fusion}
% ==============================================================================

We now test what the classical and one-loop boundary charges imply for soliton
absorption and KP-triangle defect fusion at $O(\hb^0)$
(\sfus).

\subsection{Obstruction 1: No elementary single-soliton absorption connects two stable vacua}

\begin{proposition}[Complete growing-mode count on KP triples, \Lc, \sfus]\label{prop:stability_obstruction}
Parametrise any nondegenerate KP triple $\{q_1,q_2,q_3\}$ ($q_1+q_2+q_3=0$,
$q_1 q_2 + q_2 q_3 + q_3 q_1 = -\frac{1}{4}$) by
$q_j = -\frac{1}{\sqrt{3}}\cos(\vartheta + \frac{2\pi j}{3})$ with
$\vartheta \in (0,\pi)\setminus\{\frac{\pi}{3},\frac{2\pi}{3}\}$. For either
defect orientation $\eta \in \{+1,-1\}$, the three effective defect parameters
$s_{\mathrm{eff},j} = -2\sqrt{3}\,\eta\,q_j = 2\eta\cos(\vartheta+\frac{2\pi j}{3})$
lie in $(-2,-1)$, $(-1,1)$, and $(1,2)$ respectively. By Proposition~\ref{prop:det_branches}(3):
\begin{enumerate}
\item The sorted number of strictly positive admissible growing modes ($S>1$) on
the three vertices $\{q_1,q_2,q_3\}$ is $(0,1,2)$ for all
$\vartheta \not\equiv \frac{\pi}{6} \pmod{\frac{\pi}{3}}$, and $(0,1,1)$ at the
isolated branch-switch angles $\vartheta \equiv \frac{\pi}{6} \pmod{\frac{\pi}{3}}$
where the $(1,2)$ vertex hits $s_{\mathrm{eff}} = \sqrt{3}$ ($g_2(\sqrt{3})=0$ at
threshold $S=1$, while $g_1(\sqrt{3}) = \frac{2\sqrt{3}-1}{8} \approx 0.3080 > 0$).
\item In every nondegenerate KP triple (checked across $594$ directed edges),
\textbf{exactly one vertex $q_{\mathrm{st}}$ is linearly stable} ($s_{\mathrm{eff}} < -1$),
and \textbf{both excited vertices are linearly unstable}. Consequently, \textbf{no}
elementary single-soliton absorption or emission process
$D(q_i) + S_i \leftrightarrow D(q_{i+1})$ connects two linearly stable external
defect states.
\end{enumerate}
\end{proposition}

\subsection{Obstruction 2: Non-zero KP-triangle holonomy of canonical collective-coordinate charges}

Even if one continues the one-loop defect vacuum charges
$(\Delta E_D^{\mathrm{ren}}(q_i,\eta),\Delta P_D(q_i,\eta))$ analytically from
the stable region $s_{\mathrm{eff}}<-1$ to the two unstable vertices of a KP
triple, a second obstruction arises from the bulk soliton velocity anomaly
(\texttt{charge\_frame\_master.pdf}, \texttt{audit\_charges\_action.pdf}). Recall
that for a bulk soliton $S = (\alpha,\epsilon)$, the fixed-profile mode-sum
charges $(\Delta E_1^{\mathrm{ren}}(\alpha),\Delta P_{\mathrm{mode}}(\alpha,\epsilon))$
violate $\dd E = v\,\dd P$ by $\mathcal{A}_v(\alpha) = \frac{\dd\Delta E_1^{\mathrm{ren}}}{\dd\alpha} - v\frac{\dd\Delta P_{\mathrm{mode}}}{\dd\alpha} \ne 0$,
whereas the co-moving canonical collective-coordinate derivatives
$\Delta P_{\mathrm{can}} = \frac{\partial L_1}{\partial v}$ and
$\Delta E_{\mathrm{can}} = v\Delta P_{\mathrm{can}} - L_1$ ($L_1 = v\Delta P_{\mathrm{mode}}-\Delta E_1^{\mathrm{ren}}$)
satisfy $\dd\Delta E_{\mathrm{can}} = v\,\dd\Delta P_{\mathrm{can}}$ \Lc.

\begin{proposition}[KP-triangle holonomy no-go theorem for canonical soliton charges, \Lc, \sfus]\label{prop:holonomy_nogo}
Let $\{q_1,q_2,q_3\}$ be a nondegenerate KP triple with edges $S_i = \{q_i,q_{i+1}\}$
($i=1,2,3$, $q_4\equiv q_1$) and absorption signs $\rho_i = \eta\sgn(q_i-q_{i+1})$.
\begin{enumerate}
\item Under minimal tadpole subtraction ($\delta s_1 = 0$), the 2-root dipole
defect shifts match the \textbf{mode-sum} soliton charges on every edge,
$\Delta(E_D^{\mathrm{ren}},P_D)_{q_i\to q_{i+1}} = \rho_i(\Delta E_1^{\mathrm{ren}}(S_i),\Delta P_{\mathrm{mode}}(S_i))$,
and their directed sum around the closed KP triangle vanishes identically
(verified to $8.02\times 10^{-51}$):
\begin{equation}\label{eq:mode_holonomy_zero}
  \sum_{i=1}^3 \rho_i\,\Delta E_1^{\mathrm{ren}}(S_i) = 0, \qquad
  \sum_{i=1}^3 \rho_i\,\Delta P_{\mathrm{mode}}(S_i) = 0.
\end{equation}
\item By contrast, the co-moving canonical collective-coordinate charges
$(\Delta E_{\mathrm{can}},\Delta P_{\mathrm{can}})$ have a \textbf{non-zero
directed holonomy around every generic KP triangle}:
\begin{equation}\label{eq:can_holonomy_nonzero}
  \boxed{\;
  \mathcal{C}_E(q_1,q_2,q_3) := \sum_{i=1}^3 \rho_i\,\Delta E_{\mathrm{can}}(S_i) \ne 0, \qquad
  \mathcal{C}_P(q_1,q_2,q_3) := \sum_{i=1}^3 \rho_i\,\Delta P_{\mathrm{can}}(S_i) \ne 0.
  \;}
\end{equation}
For example, for $\eta=+1$:
\begin{center}\small
\begin{tabular}{@{}ccc@{}}\toprule
KP seed root $q_1$ & $\mathcal{C}_E(q_1,q_2,q_3)$ & $\mathcal{C}_P(q_1,q_2,q_3)$ \\\midrule
$-0.20$ & $+0.3175373286$ & $+0.1655848943$ \\
$+0.10$ & $+0.4898832009$ & $-0.1111527059$ \\
$+0.45$ & $-0.5020709674$ & $+0.09965754968$ \\\bottomrule
\end{tabular}
\end{center}
\item \textbf{No-go consequence:} Because any single-valued assignment of
renormalised defect vacuum charges $(E_D^{\mathrm{quant}}(q_i),P_D^{\mathrm{quant}}(q_i))$
on the vertices $\{q_1,q_2,q_3\}$---regardless of what finite or divergent
boundary counterterms $\delta E_D^{\mathrm{ct}}(q_i),\delta P_D^{\mathrm{ct}}(q_i)$
are added at the defect---satisfies $\sum_{i=1}^3\big(Q_D^{\mathrm{quant}}(q_{i+1})-Q_D^{\mathrm{quant}}(q_i)\big) \equiv 0$,
\textbf{no boundary counterterm or defect state-charge assignment can balance
$(\Delta E_{\mathrm{can}}(S_i),\Delta P_{\mathrm{can}}(S_i))$ at the unshifted
classical absorption rapidities across all three edges of a KP triangle.}
\item Only if each absorption rapidity shifts at $O(\hb)$ relative to the defect
vertex by the conditional kinematic shift of
\texttt{audit\_charges\_action.pdf},
\begin{equation}\label{eq:a1_shift}
  \alpha_{S_i}^{\mathrm{abs}} = \alpha_{S_i} + \hb\,a_{1,i}, \qquad
  a_{1,i} = \frac{\Delta P_{\mathrm{mode}}(S_i) - \Delta P_{\mathrm{can}}(S_i)}{P_{\mathrm{cl}}'(S_i)},
\end{equation}
do the shifted canonical soliton charges
$E_{\mathrm{cl}}(S_i)+\hb(E_{\mathrm{cl}}'(S_i)a_{1,i}+\Delta E_{\mathrm{can}}(S_i)) = E_{\mathrm{cl}}(S_i)+\hb\Delta E_1^{\mathrm{ren}}(S_i)$
and $P_{\mathrm{cl}}(S_i)+\hb(P_{\mathrm{cl}}'(S_i)a_{1,i}+\Delta P_{\mathrm{can}}(S_i)) = P_{\mathrm{cl}}(S_i)+\hb\Delta P_{\mathrm{mode}}(S_i)$
reduce back to the mode-sum charges and close around the triangle (verified to
$<10^{-49}$). Establishing \eqref{eq:a1_shift} physically would require a
one-loop quantum sewing/bound-state pole calculation \Ld.
\end{enumerate}
\end{proposition}

\subsection{Obstruction 3: Absorbed-soliton momentum renormalisation in fixed-momentum fusion}

In Phase~2g (\texttt{oneloop\_defect\_phase.pdf}, Claim~6.1), one-loop defect
fusion for a test soliton $3 = (\alpha_3,\epsilon_3)$ was written as the sum of
the continuum integral identity
\begin{equation}\label{eq:ISD_fusion}
  I_{SD}(3;q_{i+1},\eta) - I_{SD}(3;q_i,\eta) = \rho_i\,I_{SS}(3,S_i)
\end{equation}
and the test-soliton momentum-renormalisation shift
$-\big(\Delta x_D(3;q_{i+1},\eta)-\Delta x_D(3;q_i,\eta)\big)\Delta P(3) = -\rho_{\mathrm{tree},i}\Delta X_3(3,S_i)\Delta P(3)$
($\rho_{\mathrm{tree},i} = \epsilon_3\rho_i$), both of which sum to zero around
a closed KP triangle \Lc.

However, when the two-soliton S-matrix $S(P_3^{\mathrm{phys}},P_{S_i}^{\mathrm{phys}})$
is expressed at fixed physical momenta of \textbf{both} colliding solitons $3$
and $S_i$, expanding the tree-level phase $\delta_0(3,S_i)$ around
$(P_3^{\mathrm{phys}},P_{S_i}^{\mathrm{phys}})$ produces two Taylor terms:
\begin{align}\label{eq:two_taylor_terms}
  \delta_{12,\mathrm{renorm}}(3,S_i)
  &= -\frac{\partial\delta_0(3,S_i)}{\partial P_3}\Delta P(3)
    - \frac{\partial\delta_0(3,S_i)}{\partial P_{S_i}}\Delta P(S_i) \nonumber\\
  &= \Delta X_3(3,S_i)\,\Delta P(3) + \Delta X_{S_i}(3,S_i)\,\Delta P(S_i),
\end{align}
where $\Delta X_{S_i}(3,S_i) = +k_{S_i}^{-1}\ln|A_{12}(3,S_i)|$.

\begin{proposition}[Non-closure of the absorbed-soliton momentum shift around a KP triangle, \Lc, \sfus]\label{prop:absorbed_renorm}
Using $\ln|A_{12}(3,S_i)| = \rho_{\mathrm{tree},i}\big(\ln|z_t(3;q_{i+1},\eta)| - \ln|z_t(3;q_i,\eta)|\big)$
and summing by parts around any closed KP triangle $\{q_1,q_2,q_3\}$ gives the
exact identity (verified to $2.41\times 10^{-50}$):
\begin{equation}\label{eq:sbp_absorbed}
  \boxed{\;
  \sum_{i=1}^3 \rho_{\mathrm{tree},i}\,\Delta X_{S_i}(3,S_i)\,\Delta P(S_i)
  = \sum_{i=1}^3 \left(\frac{\Delta P(S_{i-1})}{k_{S_{i-1}}} - \frac{\Delta P(S_i)}{k_{S_i}}\right)\ln\big|z_t(3;q_i,\eta)\big| \ne 0.
  \;}
\end{equation}
Because $\Delta P(S_i)/k_{S_i}$ varies across the three edges $S_1,S_2,S_3$, the
right-hand side of \eqref{eq:sbp_absorbed} is a non-zero function of the test
soliton rapidity $\alpha_3$ with logarithmic singularities at $p_{3,a/b}=q_i$
(e.g.\ at $q_1=-0.2,\alpha_3=0.35,\epsilon_3=\eta=+1$, it equals $-0.7742169719$
for $\Delta P_{\mathrm{mode}}$ and $-1.737010868$ for $\Delta P_{\mathrm{can}}$).
Moreover, because $\frac{\partial}{\partial q_i}\delta_D^{\mathrm{trans}}(3;q_i,\eta)$
has simple poles $\frac{\ell_3-\ell(q_i)}{p_{3,a/b}-q_i}$ that are absent from
\eqref{eq:sbp_absorbed}, no $\alpha_3$-independent one-loop shift $q_i \mapsto q_i + \hb\,\delta q_i$
of the defect vertices can cancel \eqref{eq:sbp_absorbed} for all test solitons $3$.
\end{proposition}

\subsection{What survives at \texorpdfstring{$O(\hb^0)$}{O(hbar\^0)}: precise statement}

Combining Propositions~\ref{prop:det_branches}--\ref{prop:absorbed_renorm}, we
state precisely what holds \Lc\ and what does not survive as a physical
bootstrap selection rule \Ld:
\begin{enumerate}
\item \textbf{Well-posed one-loop transmission on the stable defect ground state \Lc:}
For any defect in its unique linearly stable vacuum $q_{\mathrm{st}}$
($s_{\mathrm{eff}} = \eta\sigma_{\mathrm{st}} < -1$), the linearised spectrum has
no growing modes, phonon scattering is reflectionless ($|z_{\mathrm{ph}}(q)|=1$),
the localised momentum operator expectation value equals $\Delta P_D(q_{\mathrm{st}},\eta)$
(Proposition~\ref{prop:PD2_identity}), the linear Casimir UV divergence is
absorbed by the boundary counterterm \eqref{eq:delta_ED_ct} up to one integrable
parameter shift $\delta s_1$ (Proposition~\ref{prop:cutoff_renorm}), and
$N$-soliton transmission factorises at one loop into pairwise phases.
\item \textbf{Surviving algebraic KP-chord identities \Lc:}
At fixed bare KP root parameters $(p_{j,a},p_{j,b},q_i)$, the continuum phase
integrals obey the exact algebraic chord difference \eqref{eq:ISD_fusion} and
closed-triangle identity $\sum_{i=1}^3\rho_i\,I_{SS}(3,S_i)=0$, together with
test-soliton delay closure $\sum_{i=1}^3\rho_{\mathrm{tree},i}\Delta X_3(3,S_i)\Delta P(3)=0$
and mode-sum charge balance \eqref{eq:mode_holonomy_zero}. Likewise, a composite
two-soliton process $D(q_{\mathrm{st}}) + S_{12} + S_{23} \to D(q_{\mathrm{st}}) + S_{13}$
begins and ends on the unique stable defect vacuum $q_{\mathrm{st}}$.
\item \textbf{Why physical KP-triangle fusion does not fix the $O(\hb^0)$ two-soliton S-matrix \Lc/\Ld:}
Promoting the algebraic triangle identities to a physical quantum bootstrap
selection rule (e.g.\ to exclude $\pm\frac{\pi}{2}$ exchange constants or
$\Gamma$-factors in $S_{12}$) is obstructed at $O(\hb^0)$ because:
(i) factoring the composite process into elementary steps traverses the two
linearly unstable defect vacua $q_2,q_3$ ($\eta\sigma_{2,3}>-1$), which are not
asymptotic quantum states;
(ii) canonical soliton charges $(\Delta E_{\mathrm{can}},\Delta P_{\mathrm{can}})$
have non-zero KP-triangle holonomy \eqref{eq:can_holonomy_nonzero} unless the
absorption locus shifts at $O(\hb)$ by \eqref{eq:a1_shift}; and
(iii) holding the absorbed solitons' physical momenta fixed introduces the
non-closing remainder \eqref{eq:sbp_absorbed}.
\end{enumerate}

% ==============================================================================
\section{Summary of claims, verification scripts, and failures recorded}
\label{sec:failures}
% ==============================================================================

\begin{table}[htbp]
\centering
\footnotesize
\begin{tabular}{@{}p{4.8cm} c p{4.6cm} p{3.8cm}@{}}
\toprule
\textbf{Claim / Formula} & \textbf{Label} & \textbf{Verification script} & \textbf{Precision / Result} \\
\midrule
Classical defect charges $(E_D,P_D,\Pi_D,M_D)$ \eqref{eq:ED_proj}, \eqref{eq:PD_proj}--\eqref{eq:MD_proj} and off-shell conservation (Prop.~\ref{prop:classical_charges}) &
\Lc & \sclass &
Exact symbolic SymPy identity \\
\addlinespace
KP-triple vacuum jump factorisation \eqref{eq:vac_jump_all4} and time-dependent conservation in transmission and absorption (Prop.~\ref{prop:classical_absorption}) &
\Lc & \sclass &
Exact symbolic + $50$-digit quadrature ($<10^{-50}$) \\
\addlinespace
Corrected $4\times 4$ sewing determinant prefactor $-\frac{2}{3}$ \eqref{eq:det_corrected}, threshold factorisation \eqref{eq:P_at_S1}, and branch switch at $s_{\mathrm{eff}}=\sqrt{3}$ \eqref{eq:g_admissible} (Prop.~\ref{prop:det_branches}) &
\Lc & \sren &
Exact symbolic + $55$-digit root check \\
\addlinespace
Pointwise operator identities $\langle P_{D,2}\rangle_{\mathrm{sym}} = \frac{q}{4\pi}\theta_{D,\mathrm{odd}}'$ \eqref{eq:PD2_pointwise} and $-\langle a_t F_1+b_t G_1\rangle_{\mathrm{sym}} = \frac{\omega}{4\pi}\theta_{D,\mathrm{even}}'$ \eqref{eq:ED2_pointwise} (Prop.~\ref{prop:PD2_identity}) &
\Lc & \sren &
Exact symbolic SymPy identity \\
\addlinespace
Cutoff Casimir asymptotics \eqref{eq:ED_dens}--\eqref{eq:ED_mode}, boundary counterterm \eqref{eq:delta_ED_ct}, and off-shell integrability constraint (Prop.~\ref{prop:cutoff_renorm}) &
\Lc & \sren &
Exact symbolic + $60$-digit check at $\Lambda=10^{22}$ ($<10^{-20}$) \\
\addlinespace
Growing-mode count $(0,1,2)$ (and $(0,1,1)$ at $s_{\mathrm{eff}}=\sqrt{3}$); zero stable-to-stable absorption edges (Prop.~\ref{prop:stability_obstruction}) &
\Lc & \sfus &
Checked across $594$ directed KP edges \\
\addlinespace
Non-zero KP-triangle holonomy of $(\Delta E_{\mathrm{can}},\Delta P_{\mathrm{can}})$ \eqref{eq:can_holonomy_nonzero} vs zero holonomy of mode-sum and $a_{1,i}$-shifted charges (Prop.~\ref{prop:holonomy_nogo}) &
\Lc & \sfus &
$50$-digit mpmath ($<10^{-49}$) \\
\addlinespace
Non-closure of absorbed-soliton momentum shift \eqref{eq:sbp_absorbed} and summation-by-parts identity (Prop.~\ref{prop:absorbed_renorm}) &
\Lc & \sfus &
$50$-digit mpmath ($<2.5\times 10^{-50}$) \\
\addlinespace
Quantum derivation of shifted absorption pole \eqref{eq:a1_shift} and unstable-resonance continuation &
\Ld & --- &
Open physical problem \\
\bottomrule
\end{tabular}
\caption{Summary of claims, labels, and verification scripts in this write-up.}
\label{tab:summary}
\end{table}

\paragraph{Failures recorded.}
In accordance with \texttt{AGENTS.md}, we record all failed checks and errors
encountered during this investigation:
\begin{enumerate}
\item \textbf{Double-subtraction of $-u_{\mathrm{CZ},x}F_{\mathrm{CZ}} - v_{\mathrm{CZ},x}G_{\mathrm{CZ}}$ in $P_D$ and unsimplified $1/\eta$ in SymPy.}
In the first version of \sclass, we defined the
total defect momentum as $\frac{1}{9}(-u_{\mathrm{CZ},x}F_{\mathrm{CZ}} - v_{\mathrm{CZ},x}G_{\mathrm{CZ}} + P_{D,\mathrm{CZ}})$
and substituted $\eta^2 \mapsto 1$ without simplifying $1/\eta \mapsto \eta$.
The off-shell check $\frac{\dd}{\dd t}(P_{\mathrm{bulk}}+P_D)=0$ failed because
CZ's $P_{D,\mathrm{CZ}}$ in \cite[(3.49)]{CZ} already incorporates
$-u_{\mathrm{CZ},x}F_{\mathrm{CZ}} - v_{\mathrm{CZ},x}G_{\mathrm{CZ}} + \tilde{P}_D$.
Using $P_D = P_{D,\mathrm{CZ}}/9$ and testing $\eta=\pm 1$ gave zero off-shell
remainder identically.
\item \textbf{Sequential dictionary substitution order in SymPy and singular transmission test point.}
In Check~2 of \sclass, a single dict call
\texttt{.subs(\{s: s\_i, pL: -w\_CZ\_f/3, eta: eta\_val\})} left $\eta$ unreplaced
inside \texttt{w\_CZ\_f} when SymPy evaluated $\eta$ before \texttt{pL};
splitting the substitutions fixed the symbolic check. In Check~3(a), our first
numerical test point $(s=-1.5,\alpha=0.55,\epsilon=+1,\eta=+1)$ had
$q_D = \frac{\sqrt{3}}{4} \in (p_b,p_a)$, so $z < 0$ (singular soliton on $x>0$)
and $\ln z$ was complex, raising a \texttt{TypeError} in \texttt{mpmath}
quadrature; choosing regular transmission points ($s=-2.5,\epsilon=+1$ and
$s=-1.5,\epsilon=-1$, where $z>0$) and using $\tanh/\operatorname{sech}$ jet
evaluations to avoid large-argument overflow gave conservation to $<10^{-50}$.
\item \textbf{Zero-endpoint vs asymptotic branch in $B_D(\Lambda)$ and sign of $\sigma_0(a-b)/2$ in $F_1,G_1$.}
In the first run of \sren, evaluating the cutoff
endpoint $B_D(\Lambda)$ with $\theta_D^{\mathrm{phys}}$ ($\theta_D^{\mathrm{phys}}(0^\pm)=0$)
failed because $\theta_D^{\mathrm{phys}}(\pm\infty) = \pm\pi\eta$ for $s<-1$;
using the asymptotic branch $\theta_D^{\mathrm{asympt}}(\pm\infty)=0$ (where the
$q=0^\pm$ jump is killed by $\omega(0)=0$) and accounting for the $+\frac{S_D}{2\pi}$
rational term in $\Delta E_D^{\mathrm{ren}}$ gave agreement to $<10^{-20}$. In
Part~3, writing $+\frac{\sigma_0}{2}(a-b)$ instead of $-\frac{\sigma_0}{2}(a-b)$
in $F_1,G_1$ broke quadratic momentum conservation $\frac{\dd}{\dd t}(P_{\mathrm{bulk},2}+P_{D,2})=0$;
fixing the sign from $\delta(\frac{\sigma^2}{12\mu}) = -\frac{\sigma_0}{2}(a-b)$
and correcting the $-q$ derivative sign in $\theta_D'(-q)$ made all quadratic
identities pass identically.
\item \textbf{Threshold/branch-switch mode counts in \sfus.}
In Check~1 of \sfus, we initially excluded
$n=67$ instead of $n=68$ ($2\times 102/3 = 68$) on the grid $\vartheta = \pi n/102$
and expected growing-mode counts $(0,1,2)$ at all nondegenerate angles. The
script returned $\{(0,1,2),(0,0,2),(0,1,1)\}$, exposing both the typo $67\to 68$
(where $\vartheta=2\pi/3$ is a degenerate threshold triple) and the genuine physical
branch-switch point at $n \in \{17,51,85\}$ ($\vartheta \equiv \frac{\pi}{6}\pmod{\frac{\pi}{3}}$,
$s_{\mathrm{eff}}=\sqrt{3}$), where the second growing mode touches $g_2(\sqrt{3})=0$
($S=1$) while $g_1(\sqrt{3}) = \frac{2\sqrt{3}-1}{8} > 0$ remains strictly unstable,
giving mode count $(0,1,1)$ and confirming zero stable-to-stable edges everywhere.
\end{enumerate}

\begin{thebibliography}{9}\small
\bibitem{CZ}
E.~Corrigan and C.~Zambon,
\emph{Adding integrable defects to the Boussinesq equation},
J.~Phys.~A \textbf{56} (2023) 405202,
\href{https://arxiv.org/abs/2306.03772}{arXiv:2306.03772}.
[arXiv abstract read online for this task; bulk/defect equations and exact
solutions rederived and verified symbolically from the action in
\texttt{corrigan-zambon/defect\_check.py},
\texttt{corrigan-zambon/verify\_defect\_solutions.py}, and
\sclass; sections 2--3 and 5.1 previously checked in
\texttt{audit\_charges\_action.pdf}.]

\bibitem{FK}
L.~D.~Faddeev and V.~E.~Korepin,
\emph{Quantum theory of solitons},
Phys.~Rep.\ \textbf{42} (1978) 1--87,
\href{https://doi.org/10.1016/0370-1573(78)90058-3}{doi:10.1016/0370-1573(78)90058-3}.
[Relevant equal-mode-count passages in Appendix~6 (eqs.~A.6.11--A.6.14) verified
in \texttt{audit\_charges\_action.pdf}.]
\end{thebibliography}

\end{document}
